% 第9章：有限比特状态无法区分标签串（定理 thm:continual-memory-example 的计数，T=3, B=2）。
% 固定随机数后，编码为确定映射：此处示例为丢弃末位 b_1b_2b_3 -> b_1b_2；
% 解码 b_1b_2 -> 分类函数 (b_1,b_2,0)。2^B=4 个状态至多恢复 4 个标签串，8 个中成功 4 个，
% 平均成功率 4/8 = 2^{B-T}。任何其他编码同样至多恢复 4 个。
\begingroup
\definecolor{FigureGreen}{HTML}{829C8B}
\definecolor{FigureGreenFill}{HTML}{E8EEE9}
\begin{tikzpicture}[figure picture,x=1mm,y=1mm]
\node[figure heading] at (6,50) {标签串 $(b_1,b_2,b_3)$};
\node[figure heading] at (50,50) {2比特状态};
\node[figure heading] at (92,50) {解码出的分类函数};
\foreach \i/\s/\ok in {0/000/1,1/001/0,2/010/1,3/011/0,4/100/1,5/101/0,6/110/1,7/111/0}{%
  \pgfmathsetmacro{\yy}{44-\i*5.6}
  \ifnum\ok=1
    \node[draw=FigureGreen,fill=FigureGreenFill,rounded corners=1mm,minimum width=12mm,minimum height=4.4mm,inner sep=0pt,font=\figurefont\footnotesize] (s\i) at (6,\yy) {\s};
  \else
    \node[draw=FigureRed,fill=FigureRedFill,rounded corners=1mm,minimum width=12mm,minimum height=4.4mm,inner sep=0pt,font=\figurefont\footnotesize] (s\i) at (6,\yy) {\s};
  \fi}
\foreach \j/\m/\d in {0/00/000,1/01/010,2/10/100,3/11/110}{%
  \pgfmathsetmacro{\yy}{41.2-\j*11.2}
  \node[draw=FigureBlue,fill=FigureBlueFill,rounded corners=1mm,minimum width=11mm,minimum height=6mm,inner sep=0pt,font=\figurefont\footnotesize] (m\j) at (50,\yy) {\m};
  \node[draw=FigureStroke,fill=white,rounded corners=1mm,minimum width=16mm,minimum height=6mm,inner sep=0pt,font=\figurefont\footnotesize] (d\j) at (92,\yy) {\d};
  \draw[figure flow] (m\j) -- (d\j);}
\foreach \i/\j in {0/0,1/0,2/1,3/1,4/2,5/2,6/3,7/3}{\draw[figure link,-{Stealth[length=1.5mm]}] (s\i.east) -- (m\j.west);}
\node[figure note,align=left,anchor=north west] at (-2,-1) {绿：解码恰为原标签串；红：与另一串共用状态，解码错误。\\每个状态只能解码出一个函数，8 个串中至多 $2^2=4$ 个被恢复。};
\end{tikzpicture}
\endgroup
